\begin{lemma}\label{l:graded-trivial-zero-nilpotent}
    %
    Если градуированная конечномерная ассоциативная алгебра $A$ имеет тривиальную
    нейтральную компоненту, то $A$ нильпотентна.
\end{lemma}
\begin{proof}
    Так как алгебра $A$ конечномерна и ее градуировка
    \[
        \Gamma \colon A = \bigoplus_{g \in G}A_g
    \]
    является прямой суммой подпространств $A$, то $s = \left|\Supp \Gamma\right| <
        \infty$.

    Возьмем $m$ однородных элементов $a_1,\dots,a_m$, где $a_i\in A_{g_i}$ и
    $g_i\in \Supp \Gamma$. Рассмотрим частичные произведения
    \[
        P_k:=a_1a_2\cdots a_k \quad (k=1,\dots,m).
    \]
    По свойству градуировки имеем $P_k\in A_{g_1g_2\cdots g_k}$. Предположим, что
    $P_m\neq 0$. Тогда $P_k\neq 0$ для всех $k\le m$ и, следовательно,
    \[
        g_1g_2\cdots g_k \in \Supp \Gamma\quad (k=1,\dots,m).
    \]
    Из условия $A_{1_G} = 0$ следует $1_G \not\in \Supp \Gamma$. Положим
    \[
        p_0\coloneq 1_G,\
        p_k\coloneq g_1g_2\cdots g_k \quad (k=1,\dots,m).
    \]
    Тогда $p_0=1_G$, а $p_1,\dots,p_m\in \Supp \Gamma$, значит все
    $p_0,p_1,\dots,p_m$ лежат в множестве $\{1_G\}\cup \Supp \Gamma$ мощности $s +
        1$.

    Если $m \geq s + 1$, то по принципу Дирихле найдутся $0\le i<j\le m$ такие, что
    $p_i=p_j$. Тогда
    \[
        g_{i+1}g_{i+2}\cdots g_j \;=\; p_i^{-1}p_j \;=\; 1_G.
    \]
    Следовательно, $ a_{i+1}a_{i+2}\cdots a_j \in A_{g_{i+1}\cdots g_j}=A_{1_G}=0,
    $ то есть $a_{i+1}\cdots a_j=0$. По ассоциативности отсюда получаем $a_1\cdots
        a_m=0$, что противоречит $P_m\neq 0$.

    \medskip

    Итак, любое произведение $m\ge s+1$ однородных элементов равно нулю. Значит,
    \[
        A^{\,s+1}=0,
    \]
    поскольку каждый элемент $A$ раскладывается по однородным компонентам, и при
    раскрытии скобок в произведении $s+1$ произвольных элементов каждое слагаемое
    является произведением $s+1$ однородных элементов и потому равно нулю.
    Следовательно, $A$ нильпотентна.
\end{proof}
